\documentclass[11pt]{article}
\usepackage[margin=1in]{geometry}
\usepackage{amsmath,amssymb,amsthm,mathtools}
\usepackage[T1]{fontenc}
\usepackage{lmodern}
\usepackage{microtype}
\usepackage[colorlinks=true,urlcolor=blue,citecolor=blue,linkcolor=blue]{hyperref}
\setlength{\parskip}{3pt}
\setlength{\parindent}{1.2em}
\newtheorem{theorem}{Theorem}
\newtheorem{lemma}{Lemma}
\theoremstyle{remark}\newtheorem*{remark}{Remark}
\newcommand{\e}{\mathrm e}
\newcommand{\T}{\mathbb T}
\newcommand{\eps}{\varepsilon}
\newcommand{\M}{\mathfrak M}
\newcommand{\m}{\mathfrak m}
\newcommand{\SSS}{\mathfrak S}
\DeclareMathOperator{\meas}{meas}
\title{A compact classical proof of \(G(4)=16\)\\
\large Sparse four-power sets and eight ordinary variables}
\author{Research draft}
\date{22 November 2026}
\begin{document}
\maketitle
\begin{abstract}
We give a compact reconstruction of the classical density route to
Davenport's theorem \(G(4)=16\). An elementary diminishing-ranges argument
of Davenport and Erd\H os supplies sufficiently large sets of sums of four
fourth powers. Two such sets, combined with eight ordinary variables, are
handled by Weyl's inequality and the circle method. A smooth ordinary
cutoff makes the major-arc calculation direct. We include the complete-sum
tail and all local positivity arguments. Davenport and Erd\H os explicitly
identified this route in 1939; no new theorem or priority for the strategy
is claimed. No Vinogradov mean-value theorem or smooth-number mean-value
estimate is used.
\end{abstract}

\section{The density input}

Let \(G(k)\) be the least integer \(s\) such that every sufficiently large
positive integer is a sum of at most \(s\) nonnegative integral \(k\)th powers.
We shall prove Davenport's classical result \cite{Davenport}.
\begin{theorem}\label{thm:main}
One has \(G(4)=16\).
\end{theorem}

The following is the quartic specialization of Davenport--Erd\H os
\cite[Theorem~1]{DE}, with parity restrictions added by taking subsets.
Their final section explicitly states that this density estimate suffices
for Theorem~\ref{thm:main}. We give the short counting argument.

\begin{lemma}\label{lem:density}
Put
\[
 (\lambda_1,\lambda_2,\lambda_3,\lambda_4)
 =\left(1,\frac{48}{55},\frac{36}{55},\frac{27}{55}\right),
 \qquad \Lambda=\sum_i\lambda_i=\frac{166}{55}.
\]
For each fixed parity pattern of four variables, the set of distinct integers
\(x_1^4+\cdots+x_4^4\), with
\(X^{\lambda_i}<x_i\le2X^{\lambda_i}\) and the prescribed parities,
has cardinality \(\gg_\eps X^{\Lambda-\eps}\).
\end{lemma}

\begin{proof}
Let \(E\) count the unrestricted solutions of
\(x_1^4+\cdots+x_4^4=y_1^4+\cdots+y_4^4\) in these ranges.
We prove \(E\ll_\eps X^{\Lambda+\eps}\). First consider solutions
with \(x_4\ne y_4\). Since
\[
 |x_1^4-y_1^4|\ll X^{4\lambda_2},\qquad
 |x_1^4-y_1^4|\gg X^3|x_1-y_1|,
 \qquad4\lambda_2-3=\lambda_4,
\]
there are \(O(X^{1+\lambda_4})\) choices for \(x_1,y_1\).
Choose also \(x_2,x_3\). The equation places \(y_2^4\) in an
interval of length \(O(X^{4\lambda_3})\). Consecutive possible values
are spaced by \(\gg X^{3\lambda_2}=X^{4\lambda_3}\), so there are
\(O(1)\) choices for \(y_2\). The same reasoning gives \(O(1)\)
choices for \(y_3\), since \(3\lambda_3=4\lambda_4\).
The remaining equation is \(x_4^4-y_4^4=B\ne0\). The factor
\(x_4-y_4\) divides \(B\), and for a fixed nonzero difference the
fourth-power difference is strictly monotone in the smaller variable.
The elementary divisor bound therefore gives \(O_\eps(X^\eps)\)
pairs. The total is \(O_\eps(X^{\Lambda+\eps})\).

When \(x_4=y_4\), remove that variable and repeat the argument with
three variables, using \(4\lambda_2-3\le\lambda_3\).
Its diagonal case reduces to two variables. For the two-variable base,
the diagonal has \(O(X^{\lambda_1+\lambda_2})\) solutions;
off the diagonal choose \(x_2,y_2\) and use the nonzero difference
factorization for \(x_1,y_1\). This costs at most
\(X^{2\lambda_2+\eps}\le X^{\lambda_1+\lambda_2+\eps}\).
This proves the asserted energy bound in every case.

The number of tuples with any fixed parity pattern is
\(\gg X^\Lambda\), while their collision count is at most \(E\).
Cauchy--Schwarz bounds the number of their distinct sums below by
\(\gg X^{2\Lambda}/E\), proving the lemma.
\end{proof}

\section{Sets, arcs, and the minor-arc gain}

Fix a large integer \(n\), put \(P=n^{1/4}\), and apply
Lemma~\ref{lem:density} with \(X=P/8\). Let \(\mathcal A_0\)
be the distinct sums obtained with all four variables even, and
\(\mathcal A_4\) those obtained with all four odd. Then
\begin{equation}\label{eq:sets}
 |\mathcal A_h|\gg_\eps P^{3+1/55-\eps},\qquad
 a\equiv h\pmod{16},\quad0<a\le P^4/64
 \quad(a\in\mathcal A_h;\ h=0,4).
\end{equation}
The last bound follows from \(x_i\le2X=P/4\).
Choose \(h_1,h_2\in\{0,4\}\) so that the representative of
\(n-h_1-h_2\pmod{16}\) is in \(\{1,\ldots,8\}\).
This is possible since the three intervals
\([1,8]\), \([5,12]\), and \([9,16]\) cover all residues.
Write \(\mathcal A=\mathcal A_{h_1}\),
\(\mathcal B=\mathcal A_{h_2}\), and \(K=|\mathcal A||\mathcal B|\).
For every \(a\in\mathcal A,b\in\mathcal B\), the integer
\(m=n-a-b\) satisfies
\begin{equation}\label{eq:mrange}
 31P^4/32\le m\le P^4,
 \qquad m\pmod{16}\in\{1,\ldots,8\}.
\end{equation}

Write \(\e(t)=\exp(2\pi it)\) and \(\T=\mathbb R/\mathbb Z\).
Fix the nonnegative smooth function
\[
 w(t)=\begin{cases}
 \exp\!\left(-\dfrac1{1-100(t-3/5)^2}\right),&|t-3/5|<1/10,\\
 0,&|t-3/5|\ge1/10.
 \end{cases}
\]
Thus \(w\) is supported in \([1/2,7/10]\) and positive in its interior.
Put
\[
 f(\alpha)=\sum_{x\in\mathbb Z}w(x/P)\e(\alpha x^4),\qquad
 A(\alpha)=\sum_{a\in\mathcal A}\e(\alpha a),\quad
 B(\alpha)=\sum_{b\in\mathcal B}\e(\alpha b).
\]
The integral
\begin{equation}\label{eq:count}
 R(n)=\int_\T f(\alpha)^8A(\alpha)B(\alpha)\e(-n\alpha)\,d\alpha
\end{equation}
is a sum of nonnegative weights. If it is positive, expand \(a,b\)
using their defining four-power representations to obtain sixteen fourth
powers summing to \(n\).

Take
\[
 Q=P^{109/110},\qquad
 \M=\bigcup_{q\le Q,(c,q)=1}
 \{\alpha\in\T:|q\alpha-c|\le QP^{-4}\},
 \qquad\m=\T\setminus\M.
\]
The arcs are disjoint for large \(P\), and
\(\meas\M\ll Q^2P^{-4}\).
Classical Weyl's inequality \cite[Lemma~2.4]{Vaughan}, whose needed
specialization is proved in the appendix, gives by partial summation
\begin{equation}\label{eq:weyl}
 |f(\alpha)|^8\ll_\eps P^{8+\eps}
 (q^{-1}+P^{-1}+qP^{-4})
 \quad\text{when }(c,q)=1, |\alpha-c/q|\le q^{-2}.
\end{equation}
For \(\alpha\in\m\), Dirichlet approximation supplies such a fraction
with \(Q<q\le2P^4/Q\). Hence \(\sup_\m|f|^8\ll P^{8+\eps}/Q\).
By orthogonality and Cauchy--Schwarz,
\begin{equation}\label{eq:minor}
 \left|\int_\m f^8AB\e(-n\alpha)\,d\alpha\right|
 \ll_\eps \frac{P^{8+\eps}}Q K^{1/2}
 =o(P^4K).
\end{equation}
Indeed, \(K^{1/2}\gg P^{3+1/55-\eps}\), and the quotient of the
bound by \(P^4K\) is at most
\(P^{-1/110+2\eps}\). This is the decisive density gain.

\section{Complete sums and eight-variable local densities}

Define
\[
 S(q,c)=\sum_{r\pmod q}\e(cr^4/q),\qquad
 B_8(q)=\sum_{(c,q)=1}|q^{-1}S(q,c)|^8.
\]
We require a tail bound slightly stronger than absolute convergence.
\begin{lemma}\label{lem:tails}
One has \(\sum_{q\ge1}q^{1/2}B_8(q)<\infty\).
\end{lemma}
\begin{proof}
For odd primes \(p\) and \(p\nmid c\), multiplicative characters
express \(S(p,c)\) as a sum of at most three nonprincipal Gauss sums.
Each has modulus \(\sqrt p\), by character orthogonality, so
\(|S(p,c)|\le3\sqrt p\).
For \(h\ge2\), summing over shifts by \(p^{h-1}\) cancels the
terms with \(p\nmid r\). The remaining terms give
\[
 S(p^h,c)=p^{h-1}\ (2\le h\le4),\qquad
 S(p^h,c)=p^3S(p^{h-4},c)\ (h\ge5).
\]
Writing \(h=4u+r\), we obtain the following bounds:
\[
\begin{array}{c|c|c}
 r& B_8(p^{4u+r})&p^{(4u+r)/2}B_8(p^{4u+r})\\\hline
 0\ (u\ge1)&p^{-4u}&p^{-2u}\\
 1\ (u\ge0)&3^8p^{-4u-3}&3^8p^{-2u-5/2}\\
 2\ (u\ge0)&p^{-4u-6}&p^{-2u-5}\\
 3\ (u\ge0)&p^{-4u-5}&p^{-2u-7/2}
\end{array}
\]
The local weighted sum over \(h\ge1\) is therefore \(O(p^{-2})\).

For \(p=2\) and odd \(c\), shifts by \(2^{h-3}\) cancel odd
terms when \(h\ge5\), giving
\(S(2^h,c)=8S(2^{h-4},c)\).
For \(1\le h\le4\), directly
\(S(2^h,c)=2^{h-1}(1+\e(c/2^h))\).
Thus \(B_8(2^{h+4})=B_8(2^h)/16\), and the weighted recurrence
has ratio \(1/4\). Its local series converges as well.
Chinese-remainder factorization makes \(B_8\) multiplicative.
The product of the local weighted sums converges, proving the lemma.
\end{proof}

For an integer \(m\), put
\[
 c_q(m)=q^{-8}\sum_{(c,q)=1}S(q,c)^8\e(-mc/q),\qquad
 \SSS_8(m)=\sum_{q\ge1}c_q(m).
\]
Lemma~\ref{lem:tails} gives absolute convergence uniformly in \(m\).
\begin{lemma}\label{lem:local}
There is a fixed \(c_0>0\) such that
\(\SSS_8(m)\ge c_0\) for every integer
\(m\pmod{16}\in\{1,\ldots,8\}\).
\end{lemma}
\begin{proof}
Let \(N_h(m;p)\) count solutions of
\(x_1^4+\cdots+x_8^4\equiv m\pmod{p^h}\).
Orthogonality grouped by character conductor, and multiplicativity, give
\[
 p^{-7h}N_h(m;p)=\sum_{j=0}^h c_{p^j}(m),\qquad
 \SSS_8(m)=\prod_p\sigma_p(m),\quad
 \sigma_p(m)=\lim_{h\to\infty}p^{-7h}N_h(m;p).
\]
We first produce a solution modulo every odd prime with a nonzero coordinate.
For \(p=3,5,7\), use between one and \(p\) ones and pad with zeros.
For \(p\ge19\), the Gauss-sum estimate gives
\[
 N_1(m;p)\ge p^7-3^8p^4>1,
 \qquad19^3>3^8.
\]
The zero vector is the only possible solution with all coordinates zero,
so a nonzero one exists even when \(m\equiv0\pmod p\).

For completeness, the three remaining primes need only a small residue check.
Let \(C_p=\{x^4\pmod p\}\), including zero.
For \(p=11\), \(|C_p|=6\), so \(C_p+C_p=\mathbb F_p\) by
pigeonhole. For \(p=13,17\), respectively,
\[
 C_{13}=\{0,1,3,9\},\quad |C_{13}+C_{13}|=10;
 \qquad C_{17}=\{0,\pm1,\pm4\},\quad |C_{17}+C_{17}|=13.
\]
Thus \(3C_p=\mathbb F_p\) in both cases, again by pigeonhole.
Represent \(m-1\) in these sumsets, add a fourth power~1, and pad to
eight coordinates. In every odd-prime case, the derivative \(4x_i^3\)
at a nonzero coordinate is a unit. Hensel lifting gives
\(\sigma_p(m)\ge p^{-7}>0\).

At the prime~2, every integer congruent to~1 modulo~16 has an odd fourth
root modulo \(2^h\) for every \(h\ge4\). The induction follows from
\[
 (x+2^{h-2})^4\equiv x^4+2^h x^3\pmod{2^{h+1}}
 \qquad(h\ge4, x\text{ odd}),
\]
which toggles the next bit. If \(m\equiv j\pmod{16}\) with
\(1\le j\le8\), distinguish one coordinate required only to be odd.
Prescribe \(j-1\) of the other seven coordinates to be~1 modulo~16 and
the remaining \(8-j\) to be~0. Every choice of these seven coordinates
with their prescribed residues can be completed by an odd fourth root.
Consequently \(N_h(m;2)\ge2^{7(h-4)}\), and
\(\sigma_2(m)\ge2^{-28}\).

Finally, the unweighted prime-power bounds in the preceding lemma give
\(\sigma_p(m)=1+O(p^{-3})\) for odd \(p\), uniformly in \(m\).
The product over all sufficiently large primes is uniformly bounded below;
the finitely many remaining factors have the lower bounds just proved.
\end{proof}

\section{Major arcs with a smooth ordinary cutoff}

Put
\[
 v(\beta)=\int_{\mathbb R}w(x/P)\e(\beta x^4)\,dx,
 \qquad V(c/q+\beta)=q^{-1}S(q,c)v(\beta)
 \quad(c/q+\beta\in\M).
\]
Since \(w\) is smooth and supported away from zero,
\begin{equation}\label{eq:vdecay}
 |v(\beta)|\ll_{L,w}P(1+P^4|\beta|)^{-L}
 \qquad\text{for every fixed }L>0.
\end{equation}
Poisson summation in residue classes modulo~\(q\) gives the uniform
approximation
\begin{equation}\label{eq:poisson}
 f(\alpha)=V(\alpha)+O_{L,w}\bigl(P(Q/P)^L\bigr)
 \qquad(\alpha\in\M).
\end{equation}
Here are the details relevant to its uniformity. The Poisson expansion is
\[
 \frac1q\sum_{\ell\in\mathbb Z}S(q,c,\ell)P
 \int_{\mathbb R}w(t)\e(\beta P^4t^4-\ell Pt/q)\,dt,
 \quad S(q,c,\ell)=\sum_{r\pmod q}\e((cr^4+\ell r)/q).
\]
The zero mode is \(V\). For \(\ell\ne0\), the linear part of the
phase derivative has size \(|\ell|P/q\); the quartic part is \(O(Q/q)\).
Since \(Q/P\to0\), their difference has modulus at least
\(|\ell|P/(2q)\) for large \(P\). Repeated integration by parts,
with normalized higher derivatives uniformly bounded, and the trivial
bound \(|S(q,c,\ell)|\le q\), give~\eqref{eq:poisson}.

Take \(L=220\). The pointwise error is \(O_w(P^{-1})\).
Since \(|f|+|V|\ll_w P\), \(|AB|\le K\), and
\(\meas\M\ll Q^2P^{-4}\), replacing \(f^8\) by \(V^8\) in
the major-arc contribution to~\eqref{eq:count} costs at most
\begin{equation}\label{eq:replace}
 O_w(KP^2Q^2)=O_w(KP^{4-1/55})=o(KP^4).
\end{equation}

For \(m=n-a-b\), define
\[
 J(m)=\int_{\mathbb R}v(\beta)^8\e(-m\beta)\,d\beta.
\]
Expanding \(A,B\) and extending the complete-sum and real-integral
tails gives
\begin{equation}\label{eq:major}
 \int_\M V(\alpha)^8A(\alpha)B(\alpha)\e(-n\alpha)\,d\alpha
 =\sum_{a\in\mathcal A,b\in\mathcal B}\SSS_8(n-a-b)J(n-a-b)
   +O_w(KP^4Q^{-1/2}).
\end{equation}
This uses no distribution property of the sets \(\mathcal A,\mathcal B\).
For each fixed \(m\), denominators \(q>Q\) contribute at most
\(P^4\sum_{q>Q}B_8(q)\ll P^4Q^{-1/2}\).
For \(q\le Q\), extending the interval
\(|\beta|\le Q/(qP^4)\) costs, by~\eqref{eq:vdecay}, at most
\(P^4(Q/q)^{-1/2}B_8(q)\). Lemma~\ref{lem:tails} sums this bound.

Let \(h(u)=\frac14u^{-3/4}w(u^{1/4})\) for \(u>0\), and zero
otherwise. Fourier inversion and scaling show that
\[
 J(m)=P^4 C_w(m/P^4),\qquad C_w(t)=(h*h*\cdots*h)(t)
 \quad\text{with eight factors}.
\]
For each \(t\in[31/32,1]\), the point \((t/8)^{1/4}\) is strictly
inside \((1/2,7/10)\), where \(w>0\). Choose the first seven
convolution variables sufficiently near \(t/8\); the eighth is then
also in the positive range. Thus \(C_w(t)>0\), and continuity and
compactness give a fixed positive lower bound on this whole interval.
Combining~\eqref{eq:mrange} with Lemma~\ref{lem:local}, every summand
in the main term of~\eqref{eq:major} is therefore \(\gg_w P^4\).
Equations~\eqref{eq:minor}, \eqref{eq:replace}, and~\eqref{eq:major}
give \(R(n)\gg_w KP^4>0\) for sufficiently large \(n\).
This proves \(G(4)\le16\).

\section{The obstruction to fifteen variables}

For every integer \(a\ge0\), the number \(31\cdot16^a\) is not a
sum of at most fifteen nonnegative fourth powers. If \(a\ge1\),
reduction modulo~16 forces every variable to be even: a fourth power is
0 or~1 modulo~16, and there are fewer than sixteen variables.
Divide by~16 and repeat, reaching~31. The only fourth powers at most~31
are 0, 1, and~16. Without a~16 the sum is at most~15; one~16 needs fifteen
further ones, hence sixteen terms; two~16s are too large.
Thus \(G(4)\ge16\), completing Theorem~\ref{thm:main}.

\begin{remark}
The compactness of this deduction comes from proving existence with
two sparse four-power sets, rather than deriving the asymptotic formula
for the unrestricted sixteen-variable representation count. The density
argument and its application to \(G(4)=16\) are classical. The smooth
cutoff above is a convenient way of presenting the major arcs; it is
an ordinary weight and imposes no restriction on prime factors.
\end{remark}

\section*{Appendix: the quartic Weyl bound}

For an interval \(I\) of at most \(P\) consecutive positive integers,
write \(T_I=\sum_{x\in I}\e(\alpha x^4)\).
Expanding \(|T_I|^2\) by differences, then applying Cauchy--Schwarz
and repeating twice, gives
\[
 |T_I|^8\ll P^4\sum_{|h_1|,|h_2|,|h_3|<P}
 \min\{P,\|24\alpha h_1h_2h_3\|^{-1}\}.
\]
Each final inner sum is a geometric progression, since three differences
of \(\alpha x^4\) have linear coefficient \(24\alpha h_1h_2h_3\).
Zero differences contribute \(O(P^7)\). Group the others by their
nonzero product and use the elementary divisor bound to obtain
\begin{equation}\label{eq:weylappendix}
 |T_I|^8\ll_\eps P^7+P^{4+\eps}
 \sum_{1\le r\le P^3}\min\{P,\|24\alpha r\|^{-1}\}.
\end{equation}
If \(|\beta-b/d|\le C/d^2\), \((b,d)=1\), then in any block of
at most \(d/(2C+2)\) consecutive integers the points \(r\beta\)
are separated modulo~1 by at least \(1/(2d)\). Ordering their distances
from~0 and summing a harmonic progression shows that, for \(U,V\ge1\),
\[
 \sum_{r\le U}\min\{V,\|r\beta\|^{-1}\}
 \ll_C (U/d+1)\bigl(V+d\log(2d)\bigr).
\]
For \(\beta=24\alpha\), reduce \(24c/q=b/d\). Then
\(q/24\le d\le q\) and \(|\beta-b/d|\le24/d^2\).
Apply the last bound with \(U=P^3,V=P\) in~\eqref{eq:weylappendix}.
For \(q\le P^4\), absorbing logarithms into \(P^\eps\) gives
\[
 |T_I|^8\ll_\eps P^{8+\eps}
 (q^{-1}+P^{-1}+qP^{-4});
\]
for \(q>P^4\) this is trivial. Finally, summation by parts bounds
\(|f(\alpha)|\) by \(\|w'\|_1\max_{1\le U\le P}|T_{[1,U]}|\).
This proves~\eqref{eq:weyl} with a constant depending on the fixed weight.

\begin{thebibliography}{9}
\bibitem{Davenport}
H.~Davenport, \emph{On Waring's problem for fourth powers},
Ann. of Math. (2) \textbf{40} (1939), 731--747.
\bibitem{DE}
H.~Davenport and P.~Erd\H os, \emph{On sums of positive integral
\(k\)-th powers}, Ann. of Math. (2) \textbf{40} (1939), 533--536,
Theorem~1 and Section~4.
\url{https://www.renyi.hu/~p_erdos/1939-07.pdf}.
\bibitem{Vaughan}
R.~C.~Vaughan, \emph{The Hardy--Littlewood Method}, second edition,
Cambridge University Press, 1997, Lemma~2.4.
\end{thebibliography}
\end{document}
